\section{Accounting for optional work}\label{sec:effort}

OBBT inside a branch-and-bound solver is optional work: the solver remains
correct without it. A budget can limit this work on the execution where the
policy runs. The question a user cares about is different: does the solver
with the policy finish sooner than the same solver without it? This section
separates the two questions. Proposition~\ref{prop:ledger} bounds the optional
work of one execution. Proposition~\ref{prop:no-comparison} shows that such a
bound implies nothing about a comparison with an independent execution of the
unchanged solver. Appendix~\ref{sec:scheduling} gives comparison guarantees
for constructions that leave the baseline execution unchanged. These are elementary accounting
and scheduling facts, and we claim no novelty for them. Their role here is to
fix which conclusions the measurements of Section~\ref{sec:experiments} can
support.

\subsection{What must be charged}\label{sec:effort-charges}

The work of a callback-based tightening policy is not the number of auxiliary
LPs it solves. For the policy of Section~\ref{sec:algorithms-policy}, a
propagation call can involve seven kinds of work:
\begin{enumerate}
\item an admission test, evaluated at every propagation call;
\item a trigger evaluation, which reads the local box and validates the
incumbent to compute a cutoff;
\item construction of the auxiliary relaxation, with exact row arithmetic and
outward rounding;
\item exact feasibility checks of cached points for screening;
\item the directional LP solves;
\item validation of each candidate bound by directed rounding, and its
application;
\item bookkeeping.
\end{enumerate}
An LP count measures only item~5. A policy that solves fewer LPs but enters
more callbacks can spend more in total, since items~2--4 and~7 are paid per
callback. If a callback costs $c_0$ outside its LP solves and an LP with its
validation costs $c_1$, a change that saves $\Delta L$ LPs but adds $\Delta C$
callbacks saves time only if $c_1\Delta L>c_0\Delta C$.
Section~\ref{sec:experiments-work} measures this trade-off: there the non-LP
time per recorded callback was about three times the time of one LP with
validation.
The implementation charges items~2--7 to a timer for the added propagator that starts after
the admission test. Item~1 is not timed. Its cost appears only in the total
elapsed time of the run. The following accounting keeps such untimed work as
a separate term.

\subsection{A ledger for the enhanced execution}\label{sec:effort-ledger}

Consider one execution of the solver with the policy. Let $N(t)$ and $E(t)$
be the native and optional work charged up to time $t$; both are
nondecreasing, and $E(0)=0$. Optional work is spent in operations. An
operation is admitted at its start time $s$ only if
\begin{equation}\label{eq:admission}
 E(s)<b+\eta N(s),\qquad b\ge0,\ \eta\ge0 .
\end{equation}
Operations are serial: an admitted operation finishes before the next one
starts. Native work may be charged at any time.

\begin{proposition}[Ledger]\label{prop:ledger}
Suppose every increment of $E$ belongs to an admitted operation and every
admitted operation adds at most $\delta\ge0$ to $E$. Then, at every time $t$,
\begin{equation}\label{eq:ledger}
 E(t)\le b+\eta N(t)+\delta,\qquad
 N(t)+E(t)\le(1+\eta)N(t)+b+\delta .
\end{equation}
If each operation instead has an enforced upper bound $r$ on its charge, and
admission requires $E(s)+r\le b+\eta N(s)$, then $E(t)\le b+\eta N(t)$ for all
$t$. If the admission tests cost $J$ in total and this cost is not charged to
$E$, the total work is at most $(1+\eta)N(t)+b+\delta+J$.
\end{proposition}

\begin{proof}
If no operation has started by time $t$, then $E(t)=0\le b$. Otherwise let $s$
be the start of the last operation admitted at or before $t$. Because
operations are serial and every increment of $E$ belongs to an admitted
operation, only this operation charges $E$ during $(s,t]$, so
$E(t)\le E(s)+\delta$. Admission and monotonicity of $N$ give
$E(s)<b+\eta N(s)\le b+\eta N(t)$, which proves the first inequality; adding
$N(t)$ proves the second. With reservation, $E(t)\le E(s)+r\le b+\eta N(s)\le
b+\eta N(t)$. The last statement adds the uncharged cost $J$.
\end{proof}

The serial premise is necessary for the stated rule. With $b=1$, $\eta=0$, and
$\delta=1$, three concurrent operations can all be admitted while $E=0$; if
each then costs one unit, $E=3>b+\delta$. A concurrent implementation must
include outstanding reservations in its admission test. Callbacks of a
sequential solver are serial.

The constant $\delta$ must cover the whole operation, not one of its parts. A
time limit on each auxiliary LP does not bound relaxation construction,
screening, or validation, and an LP solver may itself return slightly after its
limit. The policy of Section~\ref{sec:algorithms-policy} uses $b$ equal to
$10\%$ of the time limit and $\eta=0$. Each callback receives a deadline equal
to the remaining allowance and checks it between construction blocks, cached
points, and LPs. Steps between these checks are not interrupted, so the
implementation enforces no fixed $\delta$; it records the actual overrun
instead. Section~\ref{sec:experiments-work} reports the largest recorded
propagator times. The admission test is evaluated outside the timer and
contributes to $J$; the proposition gives no bound on $J$ without a further
premise on the number of propagation calls.

\subsection{A ledger does not compare executions}\label{sec:effort-comparison}

In Proposition~\ref{prop:ledger}, $N$ is the native work on the
\emph{enhanced} execution. Let $N_0$ be the work of an independent execution
of the unchanged solver on the same input. Tightened bounds change LP
solutions, branching scores, heuristics, and node order, so the native part of
the enhanced execution follows a different trajectory. The ledger contains no
premise relating $N$ to $N_0$, and none can be derived from it.

\begin{proposition}[No comparison from a ledger]\label{prop:no-comparison}
Let $b>0$ and let $g$ be any real function of two variables. There is a
deterministic search procedure and an optional operation that satisfies the
reservation form of Proposition~\ref{prop:ledger} with $\eta=0$ and $E=b$, such
that $N+E>g(N_0,b)$.
\end{proposition}

\begin{proof}
Consider an abstract search procedure whose first decision has two legal
options, explored in the order of a score. Exploring the first option first
completes a proof after one unit of work. Exploring the second option first
spends $M$ units before the proof is complete, for instance because its subtree
must be exhausted first. The unchanged procedure scores the first option
higher, so $N_0=1$. The optional operation adds a valid redundant row at cost
$b$. Its cost is reserved at the start: the test $E+b\le b$ holds at $E=N=0$.
The row changes the scores, the second option is explored first, and $N=M$.
Then $N+E=M+b$, which exceeds $g(1,b)$ for large $M$.
\end{proof}

The example is a model of search, not a constructed solver instance, and it
predicts no particular solver's runtime. It shows which premise is missing from
a speedup argument based on a budget: a statement connecting the two
trajectories. The same caution applies to the certificates of
Sections~\ref{sec:certificates} and~\ref{sec:residual}. A certified ceiling on
remaining endpoint movement bounds the benefit of a specified operator in a
specified measure. It does not bound remaining search time. A small width
reduction can enable a decisive inference, and a large one can leave the
relevant relaxation error unchanged.

Guarantees against an unchanged baseline require a construction that
preserves its execution trace. Appendix~\ref{sec:scheduling} gives two such
constructions: interleaving bounds the work until one solver finishes, and a
capped enhanced run followed by a fresh baseline run gives an additive
bound on the extra work. Both use an ideal preemptible work model with
separate solver states, fixed random tapes, and traces that depend only on
charged work. The appendix also states the limit on a first action chosen
from indistinguishable histories.

\subsection{Consequences for evaluation}\label{sec:effort-evaluation}

The results of this section lead to three requirements, which the comparison
of Section~\ref{sec:experiments} follows. First, the cost of a policy must
include every item of Section~\ref{sec:effort-charges}, and auxiliary LP counts
must be reported separately from time. Second, a claim about speed must come
from independent complete runs of the solver with and without the policy,
charging startup, construction, and validation, and keeping unsolved runs and
failures in every summary; restricting a comparison to instances solved by
every variant hides lost solves and all differences among unsolved runs.
Third, a budget statement and a remaining-benefit certificate are
statements about one execution and about an operator, respectively. Neither
predicts the effect on search time, so neither can replace the measurement.
